\section{术语与指标速查}

\subsection{核心术语对照表}
\begin{longtable}{p{2.7cm}p{3.2cm}p{5.6cm}}
\toprule
\textbf{术语} & \textbf{中文理解} & \textbf{在本讲中的作用} \\
\midrule
\endfirsthead
\toprule
\textbf{术语} & \textbf{中文理解} & \textbf{在本讲中的作用} \\
\midrule
\endhead
API Access & 黑箱接口访问 & 允许行为层实验，但难以做机制级归因。 \\
Open Weights & 权重开放 & 允许内部分析、微调和蒸馏，但仍受既有训练配方约束。 \\
Open Source & 全链路开放 & 覆盖数据、模型、训练和评测，支持系统级可复现。 \\
Access Shapes Research & 访问塑造研究 & 课程主线：可访问资源决定可提出的问题类型。 \\
Agent Controller & 执行控制器 & LLM 不只是被调用模块，而是流程控制中枢。 \\
Memory Stream & 记忆流 & 记录历史观察与行为，为后续检索与反思提供上下文。 \\
Retrieval & 检索 & 从长记忆中提取当前任务最相关状态。 \\
Reflection & 反思 & 把历史轨迹压缩成高层策略，降低重复错误。 \\
Planning & 规划 & 将长任务拆解为可验证子目标。 \\
Tool Use & 工具调用 & 把语言决策映射为真实执行动作。 \\
Orchestration & 编排 & 组织多 agent 的角色边界、上下文边界和执行顺序。 \\
Benchmark & 基准评测 & 将“看起来更强”转化为可比较、可复验结果。 \\
Failure Taxonomy & 失败分类 & 区分目标误解、工具故障、策略退化等失效类型。 \\
Replayability & 可回放性 & 让每次失败都能被重放与诊断，是复现基础。 \\
Audit Log & 审计日志 & 为安全与责任追踪提供证据链。 \\
Distribution Shift & 分布偏移 & 评估 agent 在 OOD 场景的稳定性与泛化能力。 \\
Safety Gate & 安全闸门 & 对高风险动作设置权限确认与回滚机制。 \\
Reproducibility Package & 复现封装 & 汇总配置、脚本、环境与结果，供第三方验证。 \\
Open Definition & 开放定义 & 明确“开放”边界，避免术语混用导致协作失真。 \\
Research Throughput & 研究吞吐 & 由算力、工具链和组织流程共同决定。 \\
\bottomrule
\end{longtable}

\begin{knowledgebox}{为什么要做术语速查}
Agent 研究横跨模型、系统、评测与治理，不同团队常用词汇并不一致。  
将术语标准化可以显著降低协作误解，尤其在跨机构复现实验时效果明显。
\end{knowledgebox}

\subsection{本章小结}
术语统一是低成本高收益的工程动作。它不直接提升榜单分数，但会显著提升团队的实验可比较性与知识沉淀速度。


